\section{Future Browser67 and Paired244 Evaluation}
\label{sec:browser-future}

The independent Browser collector and the provenance-complete pairing path are implemented, but the experiments described in this section remain planned. Current Browser observations are sufficient to exercise collection, pairing, missingness, and audit behavior; they are not sufficiently complete or diverse to establish whether \browserfp improves manipulation detection. The future study must therefore treat the external browser as a new evidence source whose benign variability and incremental utility are measured rather than assumed.

\subsection{Why External-Browser Evidence Requires a Separate Study}

The in-app Web runtime and an independently launched system browser expose many similarly named Web APIs, but they do not run in the same container. They may use different packages and Chromium versions, separate cookie and storage state, distinct permission domains, different processes, viewports, launch timing, network routes, and user configuration. Some values that appear stable within a WebView can legitimately differ in the external browser, and some properties can change between the two captures even when the device and user remain unchanged. A raw field inequality between the probes is therefore an observation, not an integrity violation or an attack label.

The current Browser sidecar reflects this limitation by separating fields that can be compared literally under the frozen policy from fields considered not comparable because their semantics are container-, timing-, permission-, or network-dependent. This conservative policy is useful for provenance and quality-control auditing, but it does not by itself establish a detector. Before Browser-aware relations can be promoted, the project needs normal paired measurements spanning device families, Android versions, WebView and browser versions, orientations, window states, network conditions, and repeated sessions. Those data must reveal both the stable cross-container relationships and the legitimate divergence that a strong relation must tolerate.

\subsection{Required Re-Collection and Data Admission}

The formal Browser study will begin only after the App collection release and the attack-control release have been unified. Every baseline and active phase will then attempt an independent \browserfp capture under the same frozen App, Browser, backend, and pairing contracts. Each phase must preserve either a completed pair with reconnectable receipts, payload hashes, batch records, and pair provenance, or an explicit incomplete status and reason. Browser failure will continue to leave a valid \appfp observation in the App-only view, but that observation will not be padded or admitted to metrics that require \pairedfp.

The re-collection must also establish trustworthy grouping. Baseline and active observations from one controlled scenario, repeated sessions from one stable device or provider profile, and their associated App and Browser records must remain in the same split. The inventory should report sessions, accepted App records, Browser attempts, completed pairs, App-only records, stable profiles, and scenarios separately. This accounting is necessary to avoid presenting repeated launches as independent device coverage and to measure whether Browser missingness itself is correlated with device, provider, version, or intervention.

\subsection{Comparative Experimental Design}

The central comparison will place an \appfp-only system beside a joint \pairedfp system under the same grouped splits and verified two-state facts. An external-browser-only variant will be evaluated only for tasks whose label is meaningful from Browser evidence alone; it will not be forced into settings where the manipulation is confined to the embedded WebView and no Browser effect is expected. Raw traditional models will provide direct-feature baselines, while relation-based variants will test explicit comparisons between the in-app Web runtime and the external Browser surface.

The joint evaluation will distinguish simple concatenation from semantic fusion. Direct concatenation can show whether a conventional model extracts additional information from the 67 Browser fields, but it can also exploit high-cardinality identity or version shortcuts. The relation-based variant will instead encode declared cross-container expectations, applicability, and tolerances, such as platform-family compatibility, coarse version relationships, and fields expected to remain stable or legitimately diverge between containers. Ablations will remove the Browser fields as a whole and then remove individual Browser relation groups so that any gain can be attributed to a specific evidence family rather than to the increased dimensionality alone.

The analysis will also measure incomplete capture explicitly. Results will be reported separately for complete \pairedfp observations and for operational settings in which a fraction of Browser attempts fail. The study will compare complete-case performance, App-only fallback behavior, uncertainty or abstention under missing Browser evidence, and the distribution of failure reasons. It will not impute Browser values merely to preserve a fixed-dimensional model input.

Finally, the paired benign dataset will be used to quantify disagreement rates by external-browser and WebView version, device family, provider, orientation, and collection time. OOD evaluation will hold out device profiles, version families, providers, manipulation families, or later collection batches. This design is intended to answer whether an independent browser contributes information that survives natural cross-container variability, not merely whether it improves a metric on repeated observations from familiar devices.

\subsection{Interpretation of Possible Outcomes}

The external browser may help in several ways. It can provide a reference surface when a manipulation affects only the embedded WebView, reveal that the App Web runtime has diverged from a separately launched browser on the same device, or reduce uncertainty when the Native and WebView-host layers do not fully determine a browser-facing property. These benefits are hypotheses rather than present results.

Browser evidence can also be neutral or harmful. Independent updates, different storage and permission state, launch-time changes, and user customization may add legitimate disagreement that masks a useful signal or increases false alerts. A null aggregate gain would therefore remain informative if the study identifies which relation groups are redundant, which are too variable, and which specific manipulation families benefit from the external reference. A negative result would similarly constrain the role of external-browser collection and could justify retaining it for audit or dataset research rather than for online detection.

\subsection{Claim Boundary}

At the current stage, the defensible present-tense claim is infrastructural: HybridGuard can collect \browserfp, construct provenance-complete \pairedfp records when both surfaces succeed, and preserve valid App-only observations without imputation when Browser capture is incomplete. The manuscript does not describe \pairedfp as an evaluated 244-dimensional detector and does not claim that \browserfp improves detection. Such claims become available only after the unified re-collection, benign calibration, frozen relation policy, and stable-group-separated comparison described above have been completed.
